Целевой артефакт~--- не ETL и не GUI для защиты, а лёгкий сервис, который ставится рядом с уже существующим мониторингом.

\section{Требования}

\begin{itemize}
  \item источник метрик задаётся конфигом, а не кодом (Prometheus URL + PromQL либо Zabbix API);
  \item те же признаки, что в ETL, на коротком окне последних $N$ точек;
  \item сохранённая модель (joblib / ONNX);
  \item алерт в произвольный канал через webhook;
  \item поставка Docker-образом без GUI и без прогона миллионов строк pandas.
\end{itemize}

\section{Контур конфигурации}

Черновик контракта (реализация~--- после выбора модели):

\begin{verbatim}
source:
  type: prometheus   # или zabbix
  url: http://prometheus:9090
  scrape_interval_sec: 30
  queries: [...]
sink:
  type: webhook
  url: https://example.invalid/alert
model:
  path: /models/detector.joblib
  window: 15
\end{verbatim}

Zabbix подключается вторым коннектором в ту же внутреннюю схему из шести колонок.
Крупные рабочие выгрузки не коммитятся в публичный репозиторий: каталоги \texttt{datasets/raw/ZABBIX/} и \texttt{datasets/raw/PROMETHEUS/live/}.
